\section{The inversion formula}

A fundamental question about characteristic functions
is whether they contain all the information about a
distribution, or in other words whether knowing the
characteristic function determines the distribution
uniquely. This question is answered (affirmatively) by
the following theorem, which is a close cousin of the
standard inversion formula from analysis for the
Fourier transform.

\begin{thm}[The inversion formula]
\label{thm-inversion}
If $X$ is a r.v. with distribution $\mu_X$, then for
any $a<b$ we have

\begin{eqnarray*} \lim_{T\to\infty} \frac{1}{2\pi} \int_{-T}^T \frac{e^{-iat}-e^{-ibt}}{it}\varphi_X(t)\,dt &=& \mu_X((a,b)) + \frac12
\mu_X(\{a,b\}) \\ &=& \mathrm{P}(a<X<b) + \frac12 \mathrm{P}(X=a)+\frac12 \mathrm{P}(X=b).
\end{eqnarray*}

\end{thm}

\begin{cor}
\label{cor:charfun-uniqueness}
If $X,Y$ are r.v.s such that
$\varphi_X(t)\equiv \varphi_Y(t)$ for all $t\in\R$
then $X \stackrel{d}{=} Y$.
\end{cor}

\begin{xexercise}
Explain why \hyperref[cor:charfun-uniqueness]{Corollary~\ref*{cor:charfun-uniqueness}}
follows from the inversion formula.
\end{xexercise}

\begin{proof}[Proof of \texorpdfstring{\hyperref[thm-inversion]{Theorem~\ref*{thm-inversion}}}{Theorem}]
Throughout the proof, denote $\varphi(t)=\varphi_X(t)$
and $\mu=\mu_X$. For convenience, we use the notation
of Lebesgue-Stieltjes integration with respect to the
measure $\mu$, remembering that this really means
taking the expectation of some function of the r.v.
$X$. Denote

\begin{align*}
\tag{12}\label{eq:inversionproof}
I_T = \int_{-T}^T \frac{e^{-i a t}-e^{-i b t}}{it} \varphi(t)\,dt
= \int_{-T}^T \int_{-\infty}^\infty \frac{e^{-i a t}-e^{-i b t}}{it} e^{itx}d\mu(x)\,dt.
\end{align*}

\noindent Since
$\frac{e^{-i a t}-e^{-i b t}}{it} = \int_a^b e^{-i t
y}\,dy$
is a bounded function of $t$ (it is bounded in
absolute value by $b-a$), it follows by Fubini's
theorem that we can change the order of integration,
so

\begin{eqnarray*}
I_T&=& \int_{-\infty}^\infty \int_{-T}^T \frac{e^{-i a t}-e^{-i b t}}{it} e^{itx} dt\,d\mu(x) \\
&=&
\int_{-\infty}^\infty \left[ \int_{-T}^T \frac{\sin(t(x-a))}{t}\,dt-\int_{-T}^T \frac{\sin(t(x-b))}{t}\,dt \right]d\mu(x) \\
&=& \int_{-\infty}^\infty \left( R(x-a,T)-R(x-b,T) \right) d\mu(x),
\end{eqnarray*}

\noindent where we denote
$R(\theta, T) = \int_{-T}^T \sin(\theta t)/t\,dt$.
Note that in the notation of expectations this can be
written as
$I_T = \mathbf{E}\left( R(X-a,T)-R(X-b,T)\ \right)$.
This can be simplified somewhat; in fact, observe also
that

\[
R(\theta,T)= 2\textrm{sgn}(\theta) \int_0^{|\theta|T} \frac{\sin x}{x}\,dx = 2 \textrm{sgn}(\theta) S(|\theta| T),
\]

\noindent where we denote
$S(x) = \int_0^x \frac{\sin(u)}{u}\,du$ and
$\textrm{sgn}(\theta)$ is $1$ if $\theta>0$, $-1$ if
$\theta<0$ and $0$ if $\theta=0$. By a standard
convergence test for integrals, the improper integral
$\int_0^\infty \frac{\sin u}{u}\,du =
\lim_{x\to\infty}S(x)$
converges; denote its value by $C/4$. Thus, we have
shown that
$R(\theta,T)\to \frac12\textrm{sgn}(\theta) C$ as
$T\to \infty$, hence that

\[
R(x-a,T)-R(x-b,T) \xrightarrow[T\to\infty]{} \begin{cases}
C & a<x<b, \\
C/2 & x=a\textrm{ or }x=b, \\
0 & x<a\textrm{ or }x>b.
\end{cases}
\]

\noindent Furthermore, the function $R(x-a,T)-R(x-b,T)$ is
bounded in absolute value by $2\sup_{x\ge 0} S(x)$.
It follows that we can apply the bounded convergence
theorem in \eqref{eq:inversionproof} to get that

\begin{align*}
\tag{13}\label{eq:almostinversion}
I_T \xrightarrow[T\to\infty]{} C \mathbf{E}(1_{a<X<b}) + (C/2) \mathbf{E}(\mathbf{1}_{\{X=a\}} + \mathbf{1}_{\{X=b\}}) = C \mu((a,b)) + (C/2)\mu(\{a,b\}).
\end{align*}

\noindent This is just what we claimed, minus the fact that
$C=2\pi$. This fact is a well-known integral
evaluation from complex analysis. We can also deduce
it in a self-contained manner, by applying what we
proved to a specific measure $\mu$ and specific
values of $a$ and $b$ for which we can evaluate the
limit in \eqref{eq:inversionproof} directly. This is not
entirely easy to do, but one possibility, involving an
additional limiting argument, is outlined in the next
exercise; see also Exercise 1.7.5 on p. 35 in
[Dur2010], (Exercise 6.6, p. 470 in Appendix A.6 of
[Dur2004]) for a different approach to finding the
value of $C$.
\end{proof}

\begin{xexercise}
($\textbf{Recommended for aspiring analysts}$...) For
each $\sigma>0$, let $X_\sigma$ be a r.v. with
distribution $N(0,\sigma^2)$ and therefore with
density
$f_X(x)=(\sqrt{2\pi}\sigma)^{-1} e^{-x^2/2\sigma^2}$
and characteristic function
$\varphi_X(t) = e^{-\sigma^2 t^2/2}$. For fixed
$\sigma$, apply \hyperref[thm-inversion]{Theorem~\ref*{thm-inversion}} in its weak
form given by \eqref{eq:almostinversion} (that is, without
the knowledge of the value of $C$), with parameters
$X=X_\sigma$, $a=-1$ and $b=1$,  to deduce the
identity

\[
\frac{C}{\sqrt{2\pi}\sigma} \int_{-1}^1 e^{-x^2/2\sigma^2}\,dx = \int_{-\infty}^\infty \frac{2\sin t}{t} e^{-\sigma^2 t^2/2}\,dt.
\]

\noindent Now multiply both sides by $\sigma$ and take the
limit as $\sigma\to\infty$. For the left-hand side
this should give in the limit (why?) the value
$(2C)/\sqrt{2\pi}$. For the right-hand side this
should give $2\sqrt{2\pi}$. Justify these claims and
compare the two numbers to deduce that $C=2\pi$.
\end{xexercise}

\noindent The following theorem shows that the inversion
formula can be written as a simpler connection between
the characteristic function and the density function
of a random variable, in the case when the
characteristic function is integrable.

\begin{thm}
If
$\int_{-\infty}^\infty |\varphi_X(t)|\,dt < \infty$,
then $X$ has a bounded and continuous density
function $f_X$, and the density and characteristic
function are related by

\begin{eqnarray*}
\varphi_X(t) &=& \int_{-\infty}^\infty f_X(x) e^{itx}\,dx,  \\
f_X(x) &=& \frac{1}{2\pi} \int_{-\infty}^\infty \varphi_X(t) e^{-itx}\,dt.
\end{eqnarray*}

\end{thm}

\noindent In the lingo of Fourier analysis, this is known as
the
$\textbf{inversion formula for Fourier transforms}$.

\begin{proof}
This is a straightforward corollary of \texorpdfstring{\hyperref[thm-inversion]{Theorem~\ref*{thm-inversion}}}{Theorem}. See p. 95 in either [Dur2010] or
[Dur2004].
\end{proof}
